\section{第41届复赛第1题：高楼调谐质量阻尼器的受迫振动}
\subsection{题干}
质量 $M$、劲度系数 $k$ 的建筑物简化为受迫振子，受风力 $H\cos\omega t$；其上悬挂质量 $m$、摆长 $l$ 的调谐摆锤。求稳定振幅、最优调谐条件、风频率为无阻尼共振频率 $\sqrt{0.99}$ 倍时的质量比，以及含油阻尼 $f=-\gamma v'$ 的运动方程和作用。
\begin{figure}[htbp]\centering\begin{tikzpicture}[>=Stealth,scale=.9]
\draw[thick] (-2.2,2.6)--(2.2,2.6);\foreach \x in {-2.1,-1.8,...,2.1}{\draw (\x,2.6)--++(.18,.12);}
\draw[thick,decorate,decoration={zigzag,segment length=7pt,amplitude=2pt}] (-1.35,2.6)--(-1.35,.6);\node[left] at (-1.5,1.55){$k$};
\draw[fill=gray!18,thick] (-.9,.35) rectangle (.9,1.15);\node at (0,.75){$M$};
\draw[thick] (0,1.15)--(.78,-.35);\shade[ball color=gray!35] (.78,-.35) circle(.23);\node[right] at (1.0,-.35){$m$};
\draw[->] (0,1.15) arc[start angle=-90,end angle=-62,radius=.65] node[midway,right] {$\theta$};\node[right] at (.35,.4){$l$};
\draw[->,thick] (.9,.75)--(1.8,.75) node[right] {$H\cos\omega t$};\node[below] at (0,-.95){原题图1a：建筑物与调谐质量阻尼器};
\end{tikzpicture}
\caption{第41届复赛第1题原题图1a的TikZ重绘。}\end{figure}
\subsection{参考解答整理}
小幅度下耦合方程为 $\ddot x+l\ddot\theta+g\theta=0$、$M\ddot x+kx-mg\theta=H\cos\omega t$。令 $\omega_0^2=g/l$，稳定位移振幅为
\[A=\frac{(\omega_0^2-\omega^2)H}{(\omega_0^2-\omega^2)(k-M\omega^2)-m\omega_0^2\omega^2}.\]
无阻尼建筑物在 $\omega=\sqrt{k/M}$ 共振；为最大抑制该频率，可取 \[\boxed{l=Mg/k}.\]
当 $\omega=\sqrt{0.99k/M}$ 时，无阻尼振幅为 $100H/k$；按上述调谐条件将其降至1\%需 \[\boxed{m/M=101/9900\approx1/100}.\]
含阻尼时摆锤相对油的速度分量为 $v'_x=l\dot\theta\cos\theta$、$v'_y=-l\dot\theta\sin\theta$；答案给出完整非线性方程组。阻尼既抑制共振，也使风停后机械能持续耗散，建筑逐渐停振。
\subsection{转写依据}
题干依据 \texttt{CPHO41-Q01\_transcription\_v1.md}，答案依据第1至3页与 \texttt{CPHO41-Q01\_formula\_check\_v1.csv}。图1a按原题图重绘。本题三部分分值为31、4、10分。
